\documentclass[compress,10pt,aspectratio=169]{beamer}

\usepackage[utf8]{inputenc}
\usepackage[T1]{fontenc}
\usepackage{centernot}
\usepackage{graphicx}
\usepackage{appendix}
\usepackage{bm}
\usepackage{bbm}
\DeclareMathOperator*{\argmax}{arg\,max}
\usetheme{CambridgeUS}
\useinnertheme{circles}
\usepackage{booktabs}
\usepackage{pifont}
\usepackage{tabularx}
\usepackage{centernot}
\usepackage[labelfont=bf, font=normalsize, justification=centering, belowskip=7pt, aboveskip=7pt]{caption}
\usepackage[flushleft]{threeparttable}
\usepackage{multirow}
\newcolumntype{L}[1]{>{\raggedright\let\newline\\\arraybackslash\hspace{0pt}}m{#1}}
\newcolumntype{C}[1]{>{\centering\let\newline\\\arraybackslash\hspace{0pt}}m{#1}}
\newcolumntype{R}[1]{>{\raggedleft\let\newline\\\arraybackslash\hspace{0pt}}m{#1}}

\newcommand\emphblue[1]{\textcolor{blue}{#1}}

\AtBeginSubsection[]{
  \begin{frame}[noframenumbering,plain]
  \vfill
  \centering
  \begin{beamercolorbox}[sep=8pt,center,shadow=true,rounded=true]{title}
    \usebeamerfont{title}\insertsubsectionhead\par%
  \end{beamercolorbox}
  \vfill
  \end{frame}
}



\beamertemplatenavigationsymbolsempty




\title[Signal Extraction in Hiring]{\LARGE{Signal Extraction in Hiring}\\ \Large{Evidence from Teacher Recruiting}}
\date[Venice SI, June 2026]{Venice Summer Institute\\ \smallskip
{\small ``Managing the State: A New Look at Bureaucrats and Public Administration''}\\ \medskip
24--25 June 2026}
\author[Gianmarco Le\'on-Ciliotta (UPF)]{Matteo Bobba (TSE)\\ Tim Ederer (CMU)\\ Gianmarco Le\'on-Ciliotta (UPF)\\ Christopher Neilson (Yale)\\ Marco Nieddu (University of Cagliari)}


\begin{document}

{\setbeamertemplate{footline}{}
\begin{frame}[noframenumbering,plain]
\titlepage
\end{frame}
} %this removes the footline from title page

\section{Signal Extraction in Hiring: Evidence from Teacher Recruiting}

\begin{frame}{Motivation}
    \begin{itemize}
        \item Hiring is one of the most consequential decisions organizations make
        \smallskip
        \begin{itemize}
            \item[$\Rightarrow$] In the \emphblue{public sector}, arguably \emph{the} central personnel margin: rigid wages, rare dismissals, long careers
            \smallskip
            \item[$\Rightarrow$] Quality of public services (student learning, health, justice) is partly determined at the \emphblue{recruitment stage}
        \end{itemize}
        \medskip
        \item Recruitment decisions are made by \emphblue{line managers} who infer productivity from noisy signals (resumes, tests, interviews)
        \begin{itemize}
        \smallskip
            \item[$\Rightarrow$] Discretionary screening often goes wrong, both in the private and the public sector
        \end{itemize}
        \medskip
        \item Why does discretion fail to identify productive workers despite an information advantage?
        \begin{itemize}
        \smallskip
        \item[$\Rightarrow$] \emphblue{Information frictions}: managers cannot tell good from bad applicants apart
        \medskip
        \item[$\Rightarrow$] \emphblue{Agency frictions}: managers maximize something other than productivity
        \end{itemize}
        \medskip
        \item Separating them requires observing what managers \emphblue{know}, \emphblue{learn}, and \emphblue{choose}, benchmarked against an objective measure of productivity --- \emphblue{such data rarely exist}
    \end{itemize}
\end{frame}

\begin{frame}{This Paper}
    \begin{itemize}
        \item Combine \emphblue{rich data} on recruitment of \emphblue{permanent teachers} in Per\'u with an empirical model of discretionary hiring and learning
        \begin{itemize}
        \smallskip
        \item[$\Rightarrow$] Schools conduct \emphblue{interviews} and have \emphblue{significant leeway} in the hiring decision
        \medskip
        \item[$\Rightarrow$] Observe \emphblue{candidate ranking} before and after interviews + link to student outcomes
        \medskip
        \item[$\Rightarrow$] Benchmark choices against selection-corrected \emphblue{teacher value-added}
        \end{itemize}
        \bigskip
        \item Structural model of signal extraction in hiring
        \begin{itemize}
        \smallskip
            \item[$\Rightarrow$] Schools form \emphblue{priors} about \emphblue{teacher value-added} from observables and applications
            \medskip
            \item[$\Rightarrow$] \emphblue{Update beliefs} about value-added after observing \emphblue{signal} during interview
            \medskip
            \item[$\Rightarrow$] Allow for \emphblue{other objectives} of the schools and \emphblue{learning about unobserved traits}  of the applicants
        \end{itemize}
    \end{itemize}
\end{frame}



\begin{frame}{Preview of (Preliminary) Results}

    \begin{itemize}
    \item Substantial heterogeneity in \emphblue{preferences for} and \emphblue{learning about value added}
    \begin{itemize}
    \smallskip
        \item[$\Rightarrow$] Interview reveals \emphblue{little information}: increase in VA of \emphblue{0.01$\sigma$} relative to pre-interview
    \end{itemize}
    \bigskip
        \item Decomposition exercise
       \begin{itemize}
       \smallskip
    \item[$\Rightarrow$] Eliminating \emphblue{information frictions}: \emphblue{0.09$\sigma$} increase in value-added vis-a-vis status quo
    \medskip
    \item[$\Rightarrow$] \emphblue{Aligning school preferences} on teacher value added: \emphblue{0.07$\sigma$} increase in value-added
    \medskip
    \item[$\Rightarrow$] Eliminating \emphblue{both}: \emphblue{0.19$\sigma$} increase in value-added
\end{itemize}
        \bigskip
        \item Counterfactual screening policy: provide \emphblue{value-added estimates} to schools
        \begin{itemize}
            \smallskip
            \item[$\Rightarrow$] Achieve same gain in VA as interview + avoid bias against female candidates
        \end{itemize}
    \end{itemize}
\end{frame}

\subsection{Context and Teacher Value-Added}


\begin{frame}{Allocation Mechanism for Permanent Teachers in Peru}
    \begin{itemize}
        \item Application stage
        \begin{itemize}
        \smallskip
            \item[$\Rightarrow$] Applicants take a standardized test evaluating pedagogical skills and curricular knowledge
            \medskip
            \item[$\Rightarrow$] Applicants scoring above 60\% are eligible to apply for \emphblue{permanent positions}
            \medskip
            \item[$\Rightarrow$] Submit rank-ordered lists (ROLs) expressing their preferences over schools
        \end{itemize}
        \bigskip
        \item \emphblue{Shortlists} in each school of 10 applicants maximum
        \begin{itemize}
            \smallskip
            \item[$\Rightarrow$] Applicants ranked based on \emphblue{test scores} and assigned \emphblue{sequentially} to \emphblue{top three} ranked schools
            \bigskip
        \end{itemize}
        \item School-based evaluation
        \begin{itemize}
        \smallskip
            \item[$\Rightarrow$] School attributes points based on CV, interview, classroom evaluation
            \medskip
            \item[$\Rightarrow$] \emphblue{Final ranking} based on test score + points attributed by school
        \end{itemize}
    \end{itemize}
\end{frame}

\begin{frame}{School-based Evaluation Matters}
\begin{figure}[t!]

    \begin{minipage}[b]{0.4\linewidth}
\centering\includegraphics[width=1\linewidth]{Figures/density_cumscore.png}
        \captionof*{figure}{\footnotesize \textbf{a) Teachers' Score}}
    \end{minipage}
    \begin{minipage}[b]
{0.4\linewidth}
\centering
\includegraphics[width=1\linewidth]{Figures/rev_s0_prev_by_stage.png}
        \captionof*{figure}{\footnotesize \textbf{b) Probability of a Change in Top-Rank}}
    \end{minipage}

    \end{figure}
\begin{itemize}
        \item[$\Rightarrow$] 50\% probability of \emphblue{change in top-ranked teachers} from score-based evaluation
    \end{itemize}
\end{frame}


\begin{frame}{Measuring Teacher Effectiveness}
    \begin{itemize}
        \item Potential outcome (4th grade test score) of student $\ell$ if taught by teacher $j\in\{1,...,J\}$
        $$
        Y_{\ell j}=\alpha_j+X_{\ell}'\beta_j+Z_{\ell}'\gamma+\epsilon_{\ell j},
        $$
        \begin{itemize}
        \item[$\Rightarrow$] $\alpha_j$: average effect of teacher $j$ (ATE); $\beta_j$: returns to $X_{\ell}$ when exposed to $j$ (match effect)
        \end{itemize}
        \medskip
        \item Correct for \emphblue{selection on unobserved match quality} using observed ROLs $R_j$
        \begin{itemize}
        \smallskip
            \item[$\Rightarrow$] Multinomial logit control function à la Dubbin \& McFadden (1984), estimated by \emphblue{EM algorithm}
            \medskip
            \item[$\Rightarrow$] Get \emphblue{selection-corrected} value-added estimates $(\hat{\alpha}_j,\hat{\beta}_j)$ and shrunken estimates $\hat{\theta}_j^*$
        \end{itemize}
        \medskip
        \item Estimate a \emphblue{normal prior} on $\theta_j=(\alpha_j,\beta_j)$ depending on observables $z_j$ and preference types $\tau_j$
        $$
        \theta_j|z_j,\tau_j,\hat{\Omega}_j \sim \mathcal{N}(\mu(z_j,\tau_j),\Sigma_{\theta})
        $$
    \end{itemize}
\end{frame}


\begin{frame}{Estimates of the Prior Mean}
\begin{table}[t!]
\centering
\label{tva_estimates}
\begin{threeparttable}
\scriptsize
\renewcommand{\arraystretch}{1.3}
\begin{tabular}{L{1.5cm} *{8}{C{1cm}}}
  \hline
  \hline
 & Baseline ($\tau=1$) & Baseline ($\tau=2$) & Baseline ($\tau=3$) & $\times$ Female & $\times$ Test score & $\times$ Quechua & $\times$ Exp Pub $>0$ & $\times$ Exp Priv $>0$ \\
 & (1) & (2) & (3) & (4) & (5) & (6) & (7) & (8) \\
\hline
ATE & -0.036 (0.051) & 0.005  (0.071) & -0.122 (0.056) & 0.057  (0.031) & 0.096  (0.024) & 0.007  (0.050) & 0.028  (0.037) & 0.065 (0.030) \\
Lagged score & 0.525  (0.026) & 0.593  (0.033) & 0.539  (0.028) & -0.032 (0.016) & 0.008  (0.012) & -0.009 (0.022) & -0.012 (0.018) & 0.019 (0.015) \\
Female & -0.057 (0.032) & -0.046 (0.044) & -0.059 (0.036) & -0.045 (0.021) & -0.001 (0.014) & 0.013  (0.035) & 0.010  (0.022) & -0.003 (0.020) \\
   \hline
   \hline
\end{tabular}
\begin{tablenotes}
\scriptsize \item
\end{tablenotes}
    \end{threeparttable}
\end{table}
\begin{itemize}
    \item Female, high-scoring, and experienced teachers have larger average effectiveness (ATE)
    \begin{itemize}
    \smallskip
        \item[$\Rightarrow$] Observables explain little of the variation in match effects; unobserved types predict the ATE
    \end{itemize}
\end{itemize}
\end{frame}


\begin{frame}{Estimates of the Prior Variance}
\begin{table}[t!]
\centering
\label{tva_variance}
\begin{threeparttable}
\footnotesize
\renewcommand{\arraystretch}{1.3}
\setlength{\tabcolsep}{1pt}
\begin{tabular}{L{3cm} *{4}{C{1.5cm}}}
  \hline
  \hline
  & & \multicolumn{2}{c}{Correlation}\\
\cmidrule{3-4} & SD & Lagged Score & Female   \\
 & (1) & (2) & (3) \\
\hline
ATE & 0.383 (0.010) & -0.075 (0.053) & -0.066 (0.107) \\
Lagged Score & 0.151 (0.007) & & 0.076 (0.119) \\
Female & 0.108 (0.015) \\
   \hline
   \hline
\end{tabular}
\begin{tablenotes}
\scriptsize \item
\end{tablenotes}
    \end{threeparttable}
\end{table}
\begin{itemize}
    \item \emphblue{Large residual variance} in average effectiveness ($0.38\sigma$) and match effects ($0.15\sigma$ / $0.11\sigma$)
    \begin{itemize}
    \smallskip
        \item[$\Rightarrow$] Substantial room for learning in the school-based evaluation
    \end{itemize}
\end{itemize}
\end{frame}


\begin{frame}{Evidence of Misalignment Between School Preferences and TVA}
\begin{table}[t!]
\centering
\begin{threeparttable}
\footnotesize
\renewcommand{\arraystretch}{1.4}
\begin{tabular}{L{3.2cm} *{4}{C{2cm}}}
  \hline
  \hline
  & Random selection & \multicolumn{2}{c}{$\Delta$ with top ranked candidate}\\
 \cmidrule{3-4}&  & School: Pre-interview & \textcolor{red}{Prior TVA} \\
 & (1) & (2) & \textcolor{red}{(3)}  \\
\hline
Female & 0.802 & 0.004  & \textcolor{red}{0.065} \\
Test score (std.) & 1.653 & 0.068 & \textcolor{red}{0.158}  \\
Quechua & 0.089 & 0.009 & \textcolor{red}{0.009}\\
Exp public $>$ 0 & 0.900 & 0.046 & \textcolor{red}{0.032} \\
Exp private $>$ 0 & 0.712 & -0.027 & \textcolor{red}{0.104} \\
Internal candidate & 0.112 & 0.014 & \textcolor{red}{0.006} \\
\textcolor{red}{Shrunken VA} & \textcolor{red}{0.272} & \textcolor{red}{0.010} & \textcolor{red}{0.015} \\
   \hline
   \hline
\end{tabular}
\begin{tablenotes}
\scriptsize \item
\end{tablenotes}
    \end{threeparttable}
\end{table}
\begin{itemize}
    \item Pre-interview ranking \emphblue{poorly aligned} with a value-added--optimal ranking (Prior TVA)
    \begin{itemize}
    \smallskip
        \item[$\Rightarrow$] Too little weight on test scores, women and private experience; too much on internal candidates
        \medskip
        \item[$\Rightarrow$] Pre-interview ranking selects no higher VA than the prior ($0.010$ vs.\ $0.015\sigma$): \emphblue{not superior info} $\to$ agency frictions
    \end{itemize}
\end{itemize}
\end{frame}


\begin{frame}{Evidence of Learning about Effectiveness}
\begin{table}[t!]
\centering
\begin{threeparttable}
\footnotesize
\renewcommand{\arraystretch}{1.4}
\begin{tabular}{L{3.2cm} *{4}{C{2cm}}}
  \hline
  \hline
  & Random selection & \multicolumn{2}{c}{$\Delta$ with top ranked candidate}\\
 \cmidrule{3-4}&  & School: Pre-interview & \textcolor{red}{School: Post-interview} \\
 & (1) & (2) & \textcolor{red}{(3)}  \\
\hline
Female & 0.802 & 0.004 & \textcolor{red}{-0.025}  \\
Test score (std.) & 1.653 & 0.068 & \textcolor{red}{0.026}  \\
Quechua & 0.089 & 0.009 & \textcolor{red}{0.004} \\
Exp public $>$ 0 & 0.900 & 0.046 & \textcolor{red}{0.040} \\
Exp private $>$ 0 & 0.712 & -0.027 & \textcolor{red}{-0.015} \\
Internal candidate & 0.112 & 0.014 & \textcolor{red}{0.043} \\
\textcolor{red}{Shrunken VA} & \textcolor{red}{0.272} &  \textcolor{red}{0.010} & \textcolor{red}{0.022} \\
   \hline
   \hline
\end{tabular}
\begin{tablenotes}
\scriptsize \item
\end{tablenotes}
    \end{threeparttable}
\end{table}
\begin{itemize}
    \item Post-interview ranking \emphblue{less selected on observables} --- consistent with Bayesian learning
    \begin{itemize}
    \smallskip
       \item[$\Rightarrow$] But more likely male and internal: rankings also shift on dimensions unrelated to VA
       \medskip
       \item[$\Rightarrow$] Post-interview top candidate has \emphblue{higher value-added} ($0.022$ vs.\ $0.010\sigma$): interview is \emphblue{informative about TVA}
    \end{itemize}
\end{itemize}
\end{frame}

\subsection{A Model of Signal Extraction and Counterfactuals}



\begin{frame}{Beliefs Updating}
    \begin{itemize}
        \item Schools face \emphblue{common prior} on $q_{ij}=\alpha_j+\beta_j'X_i$ before the interview
        $$
        q_{ij}|z_j,d_j,R_j, X_i \sim \mathcal{N}(\mu_q(z_j,d_j,R_j, X_i),\sigma^2_q( X_i))
        $$
        \smallskip
        \item Schools \emphblue{receive a signal} $\nu_{ij}|\theta_j,X_i \sim \mathcal{N}(q_{ij},\exp(\gamma_i))$ at interview and \emphblue{update their beliefs}
$$
\mathbb{E}[q_{ij}|z_j,d_j,R_j, X_i, \nu_{ij}]=\mu_q(z_j,d_j,R_j)+\frac{\sigma_{q}^2( X_i)}{\sigma_{q}^2( X_i)+\exp(\gamma_i)}(\nu_{ij}+b(z_j)-\mu_q(z_j,d_j,R_j)).
$$
\begin{itemize}
    \item[$\Rightarrow$] $\gamma_i \sim \mathcal{N}(\mu_{\gamma},\sigma_{\gamma}^2)$: heterogeneous ability to extract informative signals
    \medskip
    \item[$\Rightarrow$] $b(z_j)$: biased interpretation of signal depending on observable type
\end{itemize}
    \end{itemize}
\end{frame}

\begin{frame}{School Preferences}
    \begin{itemize}
        \item Expected value of school $i$ from hiring teacher $j$ in stage $g=1,2$
$$
V_{ij}^g= \delta_i\mathbb{E}[q_{ij}+\varphi\eta_{ij}|\mathcal{S}^g_{ij}]+v(z_j)+\xi_{ij}+\varepsilon^g_{ij},
$$
    \item Capture \emphblue{heterogeneous preferences} flexibly
    \begin{itemize}
    \smallskip
    \item[$\Rightarrow$] $\delta_i \sim \log \mathcal{N}(\mu_{\delta},\sigma_{\delta}^2)$: alignment with objective of maximizing value-added
    \medskip
    \item[$\Rightarrow$] $v(z_j)$: preference for observable types (i.e. gender discrimination)
    \medskip
    \item[$\Rightarrow$] $\xi_{ij}\sim \mathcal{N}(0,\sigma_{\xi}^2)$: preference for fixed unobserved traits; $\varepsilon^g_{ij}$: EVT1 transitory shock
    \end{itemize}
    \medskip
    \item Candidate ranking \emphblue{before interview} identifies $\delta_i$ and $v(z_j)$; \emphblue{change after interview} identifies $\gamma_i,\xi_{ij},b(z_j)$
    \end{itemize}
\end{frame}


\begin{frame}{School Preferences}
 \begin{table}[t!]
\label{pref_principals}
\centering
\begin{threeparttable}
\footnotesize
\renewcommand{\arraystretch}{1.3}
\setlength{\tabcolsep}{1pt}
\begin{tabular}{L{5cm} *{5}{C{1.5cm}}}
  \hline
  \hline
 & \multicolumn{1}{c}{Mean} & \multicolumn{1}{c}{SD} \\
 & (1) & (2)  \\
\hline
Value-added weight $(\delta_i)$ & 0.847 (0.123) & 0.709 (0.128)  \\
Female & -0.030 (0.049) & - \\
Unobserved traits $(\xi_{ij})$ & - & 1.270 (0.023) \\
   \hline
   \hline
\end{tabular}
    \end{threeparttable}
\end{table}
\begin{itemize}
\medskip
    \item Strong alternative motives beyond TVA with \emphblue{substantial heterogeneity}
    \begin{itemize}
    \smallskip
        \item[$\Rightarrow$] Rank correlation between TVA and preferences post-interview: 0.46 (full info benchmark: 0.62)
        \medskip
        \item[$\Rightarrow$] No evidence of gender discrimination
    \end{itemize}
\end{itemize}
\end{frame}

\begin{frame}{Learning}
 \begin{table}[t!]
\label{learning_principals}
\centering
\begin{threeparttable}
\footnotesize
\renewcommand{\arraystretch}{1.3}
\setlength{\tabcolsep}{1pt}
\begin{tabular}{L{4cm} *{5}{C{1.5cm}}}
  \hline
  \hline
 & \multicolumn{1}{c}{Mean} & \multicolumn{1}{c}{SD} \\
 & (1) & (2)  \\
\hline
Signal precision $(\gamma_i)$ & 1.421 (0.386) & 0.649 (0.259)  \\
Bias Female & -3.184 (1.066)  & - \\
   \hline
   \hline
\end{tabular}
    \end{threeparttable}
\end{table}
    \begin{itemize}\medskip
    \item Interview reveals little information about quality on average but substantial heterogeneity
    \begin{itemize}
    \smallskip
        \item[$\Rightarrow$] Strong \emphblue{heterogeneity in ability to extract signals} about TVA
        \medskip
        \item[$\Rightarrow$] \emphblue{Biased learning} against female candidates
    \end{itemize}
    \end{itemize}
\end{frame}


\begin{frame}{Decomposing Frictions}
\begin{table}[t!]
\centering
\begin{threeparttable}
\footnotesize
\renewcommand{\arraystretch}{1.4}
\begin{tabular}{L{3cm} *{5}{C{1.5cm}}}
  \hline
  \hline
  & All Applicants & \multicolumn{4}{c}{$\Delta$ with top ranked candidate post-interview}\\
 \cmidrule{3-6}&  & Status-quo & No info frictions & No agency frictions & No info + agency frictions \\
 & (1) & (2) & (3) & (4) & (5) \\
\hline
Test Score (std.) & -0.049 & 0.036 & 0.032 & 0.206 & 0.030  \\
\textcolor{red}{Value-added} & \textcolor{red}{-0.177} &  \textcolor{red}{0.015} & \textcolor{red}{0.107} & \textcolor{red}{0.087} & \textcolor{red}{0.214} \\
   \hline
   \hline
\end{tabular}
\begin{tablenotes}
\scriptsize \item
\end{tablenotes}
    \end{threeparttable}
\end{table}
\begin{itemize}
        \item \emphblue{No info frictions} increase value-added of selected candidates by \emphblue{0.11$\sigma$}
        \medskip
    \item \emphblue{No agency frictions} increase value-added of selected candidates by \emphblue{0.09$\sigma$}
        \medskip
    \item Removing \emphblue{both frictions} yields gain of \emphblue{0.21$\sigma$}
\end{itemize}

\end{frame}


\begin{frame}{Policy Experiment: Disclose Value-Added Estimates}
\begin{table}[t!]
\centering
\begin{threeparttable}
\footnotesize
\renewcommand{\arraystretch}{1.4}
\begin{tabular}{L{3cm} *{5}{C{2cm}}}
  \hline
  \hline
  & All Applicants & \multicolumn{2}{c}{$\Delta$ with top ranked candidate}\\
 \cmidrule{3-4}&  & Status-quo post-interview & Counterfactual  \\
 & (1) & (2) & (3) \\
\hline
Test Score (std.) & -0.049 & 0.036 & 0.030\\
Female & 0.795 & -0.023 & 0.005\\
\textcolor{red}{Value-added} & \textcolor{red}{-0.177} &  \textcolor{red}{0.015} & \textcolor{red}{0.018}  \\
   \hline
   \hline
\end{tabular}
\begin{tablenotes}
\scriptsize \item
\end{tablenotes}
    \end{threeparttable}
\end{table}
\begin{itemize}
    \item Eliminate interview and instead \emphblue{provide value-added estimates} to schools
    \begin{itemize}
    \smallskip
    \item[$\Rightarrow$] Achieves slightly higher value-added gains and \emphblue{avoids gender discrimination}
    \end{itemize}
\end{itemize}
\end{frame}

\subsection{}

\begin{frame}{Conclusion}
    \begin{itemize}
        \item Leverage \emphblue{rich data + structural model} to study teacher recruiting in Peru
        \begin{itemize}
        \smallskip
            \item[$\Rightarrow$] Schools form priors about teachers' value added and update beliefs after interview
            \medskip
            \item[$\Rightarrow$] Candidate ranking before and after interview separately identify preferences from learning
        \end{itemize}
        \bigskip
        \item Quantify effect of frictions and \emphblue{evaluate counterfactual policy}
        \begin{itemize}
        \smallskip
            \item[$\Rightarrow$] Estimates reveal significant heterogeneity in preferences and little learning through interview
            \medskip
            \item[$\Rightarrow$] Without frictions schools would select teachers with 0.19$\sigma$ higher value-added
            \medskip
            \item[$\Rightarrow$] Eliminating interview and give info on value-added enhances efficiency and avoids biases
        \end{itemize}
    \end{itemize}
\end{frame}

\end{document}
